\chapter{2022年浙江卷}

\section{单选题}

\begin{problem}
设集合 \(A=\{1,2\}\),\(B=\{2,4,6\}\),则 \(A\cup B=\)(\quad)

\choices
    {\(\{2\}\)}
    {\(\{1,2\}\)}
    {\(\{2,4,6\}\)}
    {\(\{1,2,4,6\}\)}
\end{problem}

\begin{problem}
已知 \(a\),\(b\in\R\),且 \(a+3\i=(b+\i)\i\)(\(\i\) 为虚数单位),则(\quad)

\choices
    {\(a=1\),\(b=-3\)}
    {\(a=-1\),\(b=3\)}
    {\(a=-1\),\(b=-3\)}
    {\(a=1\),\(b=3\)}
\end{problem}

\begin{problem}
若实数 \(x\),\(y\) 满足约束条件
\(\begin{cases}
x-2\ge 0,\\
2x+y-7\le 0,\\
x-y-2\le 0,
\end{cases}\)
则 \(z=3x+4y\) 的最大值是(\quad)

\choices
    {\(20\)}
    {\(18\)}
    {\(13\)}
    {\(6\)}
\end{problem}

\begin{problem}
设 \(x\in\R\),则``\(\sin x=1\)''是``\(\cos x=0\)''的(\quad)

\choices
    {充分不必要条件}
    {必要不充分条件}
    {充分必要条件}
    {既不充分也不必要条件}
\end{problem}

\begin{problem}
某几何体的三视图如图所示(单位:\(\mathrm{cm}\)),则该几何体的体积(单位:\(\mathrm{cm}^3\))是(\quad)

\bitmapfigure[max width=0.34\textwidth]{img/2022/zhejiang/q5_fig1.png}

\choices
    {\(22\pi\)}
    {\(8\pi\)}
    {\(\dfrac{22\pi}{3}\)}
    {\(\dfrac{16\pi}{3}\)}
\end{problem}

\begin{problem}
为了得到函数 \(y=2\sin 3x\) 的图象,只要把函数 \(y=2\sin\left(3x+\frac{\pi}{5}\right)\) 图象上所有的点(\quad)

\choices
    {向左平移 \(\dfrac{\pi}{5}\) 个单位长度}
    {向右平移 \(\dfrac{\pi}{5}\) 个单位长度}
    {向左平移 \(\dfrac{\pi}{15}\) 个单位长度}
    {向右平移 \(\dfrac{\pi}{15}\) 个单位长度}
\end{problem}

\begin{problem}
已知 \(2^a=5\),\(\log_8 3=b\),则 \(4^{a-3b}=\)(\quad)

\choices
    {\(25\)}
    {\(5\)}
    {\(\dfrac{25}{9}\)}
    {\(\dfrac{5}{3}\)}
\end{problem}

\begin{problem}
如图,已知正三棱柱 \(ABC-A_1B_1C_1\),\(AC=AA_1\),\(E\),\(F\) 分别是棱 \(BC\),\(A_1C_1\) 上的点. 记 \(EF\) 与 \(AA_1\) 所成的角为 \(\alpha\),\(EF\) 与平面 \(ABC\) 所成的角为 \(\beta\),二面角 \(F-BC-A\) 的平面角为 \(\gamma\),则(\quad)

\bitmapfigure[max width=0.28\textwidth]{img/2022/zhejiang/q8_fig1.png}

\choices
    {\(\alpha\le \beta\le \gamma\)}
    {\(\beta\le \alpha\le \gamma\)}
    {\(\beta\le \gamma\le \alpha\)}
    {\(\alpha\le \gamma\le \beta\)}
\end{problem}

\begin{problem}
已知 \(a\),\(b\in\R\). 若对任意 \(x\in\R\),\(a|x-b|+|x-4|-|2x-5|\ge 0\),则(\quad)

\choices
    {\(a\le 1\),\(b\ge 3\)}
    {\(a\le 1\),\(b\le 3\)}
    {\(a\ge 1\),\(b\ge 3\)}
    {\(a\ge 1\),\(b\le 3\)}
\end{problem}

\begin{problem}
已知数列 \(\{a_n\}\) 满足 \(a_1=1\),\(a_{n+1}=a_n-\frac13a_n^2\)(\(n\in\N^*\)),则(\quad)

\choices
    {\(2<100a_{100}<\dfrac52\)}
    {\(\dfrac52<100a_{100}<3\)}
    {\(3<100a_{100}<\dfrac72\)}
    {\(\dfrac72<100a_{100}<4\)}
\end{problem}

\section{填空题}

\begin{problem}
我国南宋著名数学家秦九韶,发现了从三角形三边求面积的公式,他把这种方法称为``三斜求积'',它填补了我国传统数学的一个空白. 如果把这个方法写成公式,就是 \(S=\sqrt{\frac14\left[c^2a^2-\left(\frac{c^2+a^2-b^2}{2}\right)^2\right]}\),其中 \(a\),\(b\),\(c\) 是三角形的三边,\(S\) 是三角形的面积. 设某三角形的三边 \(a=\sqrt2\),\(b=\sqrt3\),\(c=2\),则该三角形的面积 \(S=\)\fillinblank{}.
\end{problem}

\begin{problem}
已知多项式 \((x+2)(x-1)^4=a_0+a_1x+a_2x^2+a_3x^3+a_4x^4+a_5x^5\),则 \(a_2=\)\fillinblank{},\(a_1+a_2+a_3+a_4+a_5=\)\fillinblank{}.
\end{problem}

\begin{problem}
若 \(3\sin\alpha-\sin\beta=\sqrt{10}\),\(\alpha+\beta=\frac{\pi}{2}\),则 \(\sin\alpha=\)\fillinblank{},\(\cos 2\beta=\)\fillinblank{}.
\end{problem}

\begin{problem}
已知函数
\(f(x)=\begin{cases}
-x^2+2,&x\le 1,\\
x+\dfrac{1}{x}-1,&x>1,
\end{cases}\)
则 \(f\left(f\left(\frac12\right)\right)=\)\fillinblank{};若当 \(x\in[a,b]\) 时,\(1\le f(x)\le 3\),则 \(b-a\) 的最大值是\fillinblank{}.
\end{problem}

\begin{problem}
现有 7 张卡片,分别写上数字 \(1\),\(2\),\(2\),\(3\),\(4\),\(5\),\(6\). 从这 7 张卡片中随机抽取 3 张,记所抽取卡片上数字的最小值为 \(\xi\),则 \(P(\xi=2)=\)\fillinblank{},\(E(\xi)=\)\fillinblank{}.
\end{problem}

\begin{problem}
已知双曲线 \(\frac{x^2}{a^2}-\frac{y^2}{b^2}=1\)(\(a>0\),\(b>0\))的左焦点为 \(F\),过 \(F\) 且斜率为 \(\frac{b}{4a}\) 的直线交双曲线于点 \(A(x_1,y_1)\),交双曲线的渐近线于点 \(B(x_2,y_2)\),且 \(x_1<0<x_2\). 若 \(|FB|=3|FA|\),则双曲线的离心率是\fillinblank{}.
\end{problem}

\begin{problem}
设点 \(P\) 在单位圆的内接正八边形 \(A_1A_2\cdots A_8\) 的边 \(A_1A_2\) 上,则 \(\overrightarrow{PA_1}^{\,2}+\overrightarrow{PA_2}^{\,2}+\cdots+\overrightarrow{PA_8}^{\,2}\) 的取值范围是\fillinblank{}.
\end{problem}

\section{解答题}

\begin{problem}
在 \(\triangle ABC\) 中,角 \(A\),\(B\),\(C\) 所对边分别为 \(a\),\(b\),\(c\). 已知 \(4a=\sqrt5\,c\),\(\cos C=\frac35\).

\begin{enumerate}
\item 求 \(\sin A\) 的值;

\item 若 \(b=11\),求 \(\triangle ABC\) 的面积.
\end{enumerate}
\end{problem}

\begin{problem}
如图,已知 \(ABCD\) 和 \(CDEF\) 都是直角梯形,\(AB\parallel DC\),\(DC\parallel EF\),\(AB=5\),\(DC=3\),\(EF=1\),\(\angle BAD=\angle CDE=60^\circ\),二面角 \(F-DC-B\) 的平面角为 \(60^\circ\). 设 \(M\),\(N\) 分别为 \(AE\),\(BC\) 的中点.

\begin{enumerate}
\item 证明:\(FN\perp AD\);

\item 求直线 \(BM\) 与平面 \(ADE\) 所成角的正弦值.
\end{enumerate}

\bitmapfigure[max width=0.34\textwidth]{img/2022/zhejiang/q19_fig1.png}
\end{problem}

\begin{problem}
已知等差数列 \(\{a_n\}\) 的首项 \(a_1=-1\),公差 \(d>1\). 记 \(\{a_n\}\) 的前 \(n\) 项和为 \(S_n\)(\(n\in\N^*\)).

\begin{enumerate}
\item 若 \(S_4-2a_2a_3+6=0\),求 \(S_n\);

\item 若对于每个 \(n\in\N^*\),存在实数 \(c_n\),使 \(a_n+c_n\),\(a_{n+1}+4c_n\),\(a_{n+2}+15c_n\) 成等比数列,求 \(d\) 的取值范围.
\end{enumerate}
\end{problem}

\begin{problem}
如图,已知椭圆 \(\frac{x^2}{12}+y^2=1\). 设 \(A\),\(B\) 是椭圆上异于 \(P(0,1)\) 的两点,且点 \(Q\left(0,\frac12\right)\) 在线段 \(AB\) 上,直线 \(PA\),\(PB\) 分别交直线 \(y=-\frac12x+3\) 于 \(C\),\(D\) 两点.

\begin{enumerate}
\item 求点 \(P\) 到椭圆上点的距离的最大值;

\item 求 \(|CD|\) 的最小值.
\end{enumerate}

\bitmapfigure[max width=0.40\textwidth]{img/2022/zhejiang/q21_fig1.png}
\end{problem}

\begin{problem}
设函数 \(f(x)=\frac{\e}{2x}+\ln x\)(\(x>0\)).

\begin{enumerate}
\item 求 \(f(x)\) 的单调区间;

\item 已知 \(a\),\(b\in\R\),曲线 \(y=f(x)\) 上不同的三点 \((x_1,f(x_1))\),\((x_2,f(x_2))\),\((x_3,f(x_3))\) 处的切线都经过点 \((a,b)\). 证明:

\begin{enumerate}
\item 若 \(a>\e\),则 \(0<b-f(a)<\frac12\left(\frac{a}{\e}-1\right)\);

\item 若 \(0<a<\e\),\(x_1<x_2<x_3\),则 \(\frac{2}{\e}+\frac{\e-a}{6\e^2}<\frac1{x_1}+\frac1{x_3}<\frac2a-\frac{\e-a}{6\e^2}\).
\end{enumerate}
\end{enumerate}

(注:\(\e=2.71828\cdots\) 是自然对数的底数)
\end{problem}

